\documentclass[11pt,a4paper]{article}
\usepackage[T1]{fontenc}
\usepackage{lmodern}
\usepackage[margin=26mm]{geometry}
\usepackage{amsmath,amssymb,amsthm,booktabs,array}
\usepackage{microtype}
\usepackage[hidelinks]{hyperref}
\hypersetup{
  pdftitle={Further OEIS proofs and a corrected conjecture},
  pdfauthor={Dev Goyal},
  pdfsubject={Version 4: AI-generated proofs and Lean formalizations}
}
\setlength{\parskip}{3pt}
\setlength{\emergencystretch}{2em}
\raggedbottom
\allowdisplaybreaks[1]
\clubpenalty=10000
\widowpenalty=10000
\newtheorem{theorem}{Theorem}[section]
\newtheorem{lemma}[theorem]{Lemma}
\newtheorem{corollary}[theorem]{Corollary}
\theoremstyle{remark}
\newtheorem{remark}[theorem]{Remark}
\newcommand{\Z}{\mathbb Z}
\newcommand{\Q}{\mathbb Q}
\newcommand{\N}{\mathbb N}
\newcommand{\F}{\mathbb F}
\newcommand{\ps}[1]{#1[[x]]}
\newcommand{\Rev}{\operatorname{Rev}}
\newcommand{\coeff}[2]{[x^{#1}]#2}
\newcommand{\E}{\mathcal E}
\title{Further OEIS proofs and a corrected conjecture\\
  \large Expanded proof package, version 4}
\author{Dev Goyal}
\date{24 October 2026}

\begin{document}
\maketitle

\begin{center}
\fbox{\begin{minipage}{0.90\linewidth}
\textbf{AI disclosure and review status.}
ChatGPT (OpenAI) generated the proofs, exposition, Lean formalizations,
verification programs, and AI reviews for this contribution in response to
Dev Goyal's requests. Dev Goyal is the sole human author and initiated and
requested the work. The source sequences and conjectures are credited to
Paul D. Hanna. \textbf{Independent human review is pending.} This paper
does not claim OEIS acceptance, peer review, or worldwide priority.
\end{minipage}}
\end{center}

\begin{abstract}
We prove six assertions labeled conjectural in five retrieved OEIS
entries: A195194, A393171, A393172, A380056, and A380060. The results
describe binary support in an infinite family of compositional inverses,
sparse congruences modulo four and eight for Catalan substitution
equations, and congruences modulo four and five for an exponential
generating function. A rational theta fixed point has residue two modulo
four exactly at positive square indices. For the related explicit theta
sum A380680, the corresponding square-index conjecture is false; we
formally verify counterexamples and prove the corrected support at twice
a square. We also formalize an elementary odd-index evenness observation
for A392210. Each implicit source is constructed and proved unique, and
all coefficient conclusions are unbounded Lean theorems. Exact finite
calculations separately corroborate the identification with the source
sequences. This paper supplements the three earlier manuscripts retained
unchanged in the cumulative archive.
\end{abstract}

\section{Scope and conventions}

All power series in the arguments are formal. The notation $G\circ H$
means $G(H(x))$, with the outer series first, and is used when $H(0)=0$.
A \emph{normalized} series has the form $x+O(x^2)$. Here $O(x^r)$ means
divisibility by $x^r$, including over rings with zero divisors. We use
$\Rev(G)$ for compositional inversion and $U^{-1}$ for the multiplicative
inverse of a unit. The set $\N$ includes zero.

Except in the A380056 row, the sequences in Table~\ref{tab:results} are
ordinary coefficients $a_n=[x^n]A$. For A380056, they are exponential
coefficients $a_n=n![x^n]A$. Every power-of-two exponent $k$ is a natural
number. The index restrictions in the table are part of the assertions.

\clearpage
\begin{table}[ht]
\centering\small
\begin{tabular}{@{}lp{0.34\linewidth}p{0.41\linewidth}@{}}
\toprule
Entry & Source specification & Proved assertion\\
\midrule
A195194 &
 $F_j=\Rev(x-(j+1)xF_{j+1})$, $A=F_0$ &
 $a_n$ odd iff $n=2^k$, $n\ge1$.\\[5pt]
A393172 & $A(x-x^2)=x+A(x)^2$ &
 For $n>1$, residue $2$ modulo $4$ exactly when $n=2^k$, and $0$ otherwise.\\[5pt]
A393171 & $B=x+B(C)^2$, $C=x+C^2$ &
 For $n>2$, residue $4$ modulo $8$ exactly when $n=2^k$ or $2^k+1$ for $k\ge1$, and $0$ otherwise.\\[5pt]
A380056 & $A=(e^x-1)/\cos(2x)$ &
 $a_n\equiv1\pmod4$ for $n\ge1$;
 $5\mid a_{4n}$ for $n\ge1$.\\[5pt]
A380060 & $A=1+2\Sigma(A)$, with $\Sigma$ defined in Section~\ref{sec:theta} &
 For $n>0$, residue $2$ modulo $4$ exactly at square indices, and $0$ otherwise.\\
\bottomrule
\end{tabular}
\caption{Six conjectural assertions in five source entries proved here.}
\label{tab:results}
\end{table}

The official entry pages were checked on 24 October 2026
\cite{oeis194,oeis171,oeis172,oeis056,oeis060}. The five entries explicitly
labeled these six statements conjectures. Exact-identifier and
defining-equation searches did not identify an earlier proof of the target
claims. This bounded search records what was checked; it does not prove
global openness, first discovery, or absence of unpublished work. The
algebraic methods used below are standard.

Two additional results have distinct status. The square-index statement
in the retrieved A380680 comment is contradicted by its own displayed
coefficients. Section~\ref{sec:theta} proves a corrected theorem, rather
than asserting that the original statement is true. The A392210 comment
describes odd-index evenness as an observed pattern; its proof in
Section~\ref{sec:odd} is a direct application of characteristic-two
squaring, not a new general principle.

The ordinary series in the first three rows and Section~\ref{sec:odd}
start at index one with coefficient one. The theta series start at index
zero with coefficient one. A380056 starts at index one; we use the natural
auxiliary value $a_0=0$ when formulating its convolution recurrence.

\section{Algebraic tools}

\begin{lemma}[Triangular construction]\label{lem:construction}
Let $K$ be a commutative ring. Suppose $\Phi$ preserves normalized series
and improves agreement by one degree:
\[
 G-H=O(x^r)\quad\Longrightarrow\quad
 \Phi(G)-\Phi(H)=O(x^{r+1})\qquad(r\ge2).
\]
Then $\Phi$ has a unique normalized fixed point. The same conclusion holds
for a family of series if the degree-$n$ equations, at every family index,
determine the degree-$n$ coefficients from strictly lower degrees.
\end{lemma}

\begin{proof}
The degree-$n$ coefficient of $\Phi(G)$ depends only on coefficients of
$G$ below degree $n$. Fix the constant and linear coefficients, and then
define every later coefficient successively. This gives a fixed point;
the same induction proves uniqueness. For a family, perform the induction
on degree simultaneously at every family index. The recurrence requires
only ring operations, so it also constructs integral solutions when
$K=\Z$. Equivalently, iterating $\Phi$ from any normalized initial series
fixes one further coefficient at each step.
\end{proof}

\begin{lemma}[Inversion and transport]\label{lem:inverse}
A normalized series over any commutative ring has a unique normalized
two-sided compositional inverse. A constant-one series has a unique
multiplicative inverse with constant coefficient one. These inverses and
valid composition commute with coefficient ring homomorphisms.
Normalized compositional inversion preserves agreement modulo $x^r$.
\end{lemma}

\begin{proof}
For normalized $G$, the degree-$n$ equation in $G(H)=x$ determines $h_n$
with multiplier one from earlier coefficients. Construct a right inverse
$J$ of $H$. Associativity gives
$G=G\circ(H\circ J)=(G\circ H)\circ J=J$, so $H\circ G=x$ as well.
The same associativity argument proves uniqueness. Multiplicative
inversion follows similarly from the coefficient recurrence for $UV=1$.

Each coefficient of a valid composition is a finite polynomial in the
coefficients of its arguments. Ring homomorphisms preserve these
polynomials and the defining equations, so uniqueness gives transport of
inversion. If $G-H=O(x^r)$, composition with any normalized series
preserves that agreement. Composing with the respective inverses then
shows $\Rev(G)-\Rev(H)=O(x^r)$.
\end{proof}

\begin{lemma}[Binary series and doubling]\label{lem:binary}
There is a unique normalized integral series $C=x+C^2$. Its reduction
modulo two is
\[
 P(x)=\sum_{k\ge0}x^{2^k}.
\]
For series $U,V$ over $\Z/4\Z$, the equality $2U=2V$ is equivalent to
equality of their reductions modulo two.
\end{lemma}

\begin{proof}
The operator $G\mapsto x+G^2$ improves agreement because
$G^2-H^2=(G-H)(G+H)$ and normalized $G+H$ has zero constant coefficient.
In characteristic two, $U(x)^2=U(x^2)$ for every formal series. Thus the
normalized equation $U=x+U^2$ forces the coefficient at one to be one,
every higher odd-index coefficient to be zero, and the coefficient at
$2m$ to equal the coefficient at $m$. This gives exactly $P$.

For scalar residues modulo four, $2u=2v$ holds precisely when $u$ and $v$
have the same parity. Apply this statement coefficientwise. In
particular, this argument does not cancel a factor of two in a ring in
which two is a zero divisor.
\end{proof}

\section{A195194: infinite nested reversion}

The title of A195194 is an infinite sequence of nested reversions
\cite{oeis194}. It is useful to construct the more general family with
arbitrary weights before specializing to $1,2,3,\ldots$.

\begin{theorem}[Canonical nested tails]\label{thm:tails}
For every sequence $(w_j)_{j\ge0}$ in a commutative ring $K$, there is a
unique family of normalized series satisfying
\[
 F_j=\Rev(x-w_jxF_{j+1})\qquad(j\ge0). \tag{1}
\]
If the finite nests are defined by $F_j^{[0]}=x$ and
\[
 F_j^{[d+1]}=\Rev(x-w_jxF_{j+1}^{[d]}),
\]
then $F_j^{[d]}-F_j=O(x^{d+2})$ for every $d,j\ge0$.
\end{theorem}

\begin{proof}
The right inverse identity in (1) is equivalent to
\[
 F_j=x+w_jF_j\,F_{j+1}(F_j). \tag{2}
\]
Both factors in the nontrivial term start in degree one. Consequently its
degree-$n$ coefficient depends only on coefficients below degree $n$ of
$F_j$ and $F_{j+1}$. Lemma~\ref{lem:construction} constructs the unique
family over $K$. No boundary condition at infinity is assumed.
Since $x-w_jxF_{j+1}$ is normalized, Lemma~\ref{lem:inverse} turns the
right inverse equation into both inverse identities, proving (1).

At depth zero both the approximation and the true tail start with $x$,
so they agree modulo $x^2$. Multiplication by $x$ improves the agreement
of their next tails by one degree, and compositional inversion preserves
that improved agreement. Induction on $d$ gives the stated bound. Thus
the finite nested expressions agree with the canonical family through
degree $d+1$; their formal infinite limit is unambiguous.
\end{proof}

\begin{theorem}[A195194 parity]\label{thm:nestedparity}
Let $w_j=j+1$, and let $A=F_0$ in Theorem~\ref{thm:tails}, over $\Z$.
For every $n\in\N$,
\[
 [x^n]A\text{ is odd}\quad\Longleftrightarrow\quad
 n=2^k\text{ for some }k\in\N.
\]
\end{theorem}

\begin{proof}
Map the family to $\F_2[[x]]$ and write a bar for reduction. The second
weight is two, so (2) gives $\overline F_1=x$. The first weight is one,
and therefore
\[
 \overline A=x+\overline A\,\overline F_1(\overline A)
             =x+\overline A^2.
\]
Lemma~\ref{lem:binary} identifies $\overline A$ with $P$. The conclusion
also holds at $n=0$, where the coefficient is zero and the index is not a
power of two.
\end{proof}

The formal implementation constructs the simultaneous family using the
coefficient ring $\N\to K$, with its shift endomorphism. Evaluation at an
index recovers (2). The theorem \texttt{source\_existence\_uniqueness}
states the general source result using both inverse identities;
\texttt{approximation\_agrees} supplies the finite-depth bound, and
\texttt{a195194\_parity} gives the infinite parity statement. All are in
\texttt{Round4.NestedReversion}.

\section{A393172 and A393171: Catalan substitution}

Let $C=x+C^2$ and put $N=x-x^2$. Then $N\circ C=x$, and
Lemma~\ref{lem:inverse} also gives $C\circ N=x$. We use $A$ for A393172
and $B$ for A393171 throughout this section, to keep the two source
series distinct \cite{oeis172,oeis171}.

\begin{lemma}[A useful square equation]\label{lem:square}
For normalized $T,P\in K[[x]]$, there is a unique normalized solution of
\[
 G=T+(G\circ P)^2.
\]
The construction commutes with coefficient ring homomorphisms.
\end{lemma}

\begin{proof}
If $G-H=O(x^r)$, then $(G-H)\circ P=O(x^r)$. The sum
$G\circ P+H\circ P$ has zero constant coefficient, so the difference
of the two squares is $O(x^{r+1})$. Normalization is preserved.
Lemma~\ref{lem:construction} gives existence and uniqueness; mapping the
equation and using uniqueness gives transport.
\end{proof}

\begin{theorem}[A393172 modulo four]\label{thm:catalanfour}
The unique normalized integral solution of
\[
 A\circ N=x+A^2 \tag{3}
\]
has, for every $n>1$,
\[
 [x^n]A\equiv
 \begin{cases}
 2\pmod4,& n=2^k\text{ for some }k\ge0,\\
 0\pmod4,& \text{otherwise}.
 \end{cases}
\]
\end{theorem}

\begin{proof}
Composing (3) with $C$ gives
$A=C+(A\circ C)^2$. Conversely, composing this equation with $N$
gives (3). Lemma~\ref{lem:square} therefore proves integral existence
and uniqueness for the literal source equation.

Over $\Z/4\Z$, the normalized candidate $D=2C-x$ satisfies
\[
 D\circ N=2x-N=x+x^2,
 \qquad D^2=x^2+4(C^2-xC)=x^2.
\]
Thus $D$ solves (3), and uniqueness gives $A=2C-x$ modulo four. For
$n>1$, the coefficient of $x$ contributes nothing and Lemma~\ref{lem:binary}
shows that twice the coefficient of $C$ is two precisely at powers of
two, and zero elsewhere.
\end{proof}

\begin{theorem}[A393171 modulo eight]\label{thm:catalaneight}
There is a unique normalized integral series
\[
 B=x+(B\circ C)^2.
\]
For every $n>2$,
\[
 [x^n]B\equiv
 \begin{cases}
 4\pmod8,& n=2^k\text{ or }2^k+1\text{ for some }k\ge1,\\
 0\pmod8,& \text{otherwise}.
 \end{cases}
\]
Moreover $B=x+A^2=A\circ N$ and $B\circ C=A$.
\end{theorem}

\begin{proof}
Existence and uniqueness follow from Lemma~\ref{lem:square} with $T=x$,
$P=C$. Define $B=x+A^2$. The fixed point equation for $A$ gives
$B\circ C=C+(A\circ C)^2=A$, so this $B$ satisfies the required source
equation. Equation (3) gives $B=A\circ N$; it also implies
$B\circ N=N+B^2$.

Theorem~\ref{thm:catalanfour} provides an integral series $K$ with
$A=2C-x+4K$. Squaring over the integers and reducing modulo eight gives
\[
 \begin{aligned}
 B&\equiv x+(2C-x)^2\\
  & =x+x^2+4(C^2-xC)\\
  & =x+x^2+4(C-x-xC)\pmod8.
 \end{aligned}\tag{4}
\]
For $n>2$, the relevant coefficient divided by four has parity
$[x^n]C+[x^{n-1}]C$. By Lemma~\ref{lem:binary}, it is one precisely
when one of $n,n-1$ is a power of two. They cannot both be powers of two
for $n>2$: consecutive powers of two can only be one and two. The exact
exceptional indices are therefore the two families stated above.
\end{proof}

The passage to (4) first establishes exact divisibility by four in the
integer series; it does not divide by four in $\Z/8\Z$. The canonical
constructions, both source equivalences, the identities relating $A$ and
$B$, and both full residue classifications are formalized in
\texttt{Round4.CatalanComposition}. The general contraction construction
is in \texttt{Round4.SquareComposition}.

\section{A380056: exponential coefficients and binomial convolution}

Here the source is the exponential generating function
\cite{oeis056}
\[
 A(x)=\sum_{n\ge0}a_n\frac{x^n}{n!}
     =\frac{e^x-1}{\cos(2x)},\qquad a_0=0. \tag{5}
\]
The integer sequence must be distinguished from the ordinary rational
coefficients of this formal series. We first establish that distinction
and the source equation exactly.

For two sequences over a commutative ring, define binomial convolution by
\[
 (f\star g)_n=\sum_{k=0}^n\binom nk f_k g_{n-k}.
\]
Write $\delta=(1,0,0,\ldots)$ and $d=(0,1,1,\ldots)$, and define
\[
 c_{2j}=(-4)^j,\qquad c_{2j+1}=0.
\]

\begin{lemma}[Integral source construction]\label{lem:egfsource}
There is a unique integer sequence $a$ satisfying $a\star c=d$. Its
exponential generating function is exactly (5).
\end{lemma}

\begin{proof}
Since $c_0=1$, set $a_0=0$ and, for $n>0$, set
\[
 a_n=1-\sum_{k=0}^{n-1}\binom nk a_k c_{n-k}.
\]
This determines integers without division and gives the stated
convolution identity. Induction proves uniqueness over any commutative
coefficient ring.

Over $\Q$, let $\E(f)=\sum_{n\ge0}f_nx^n/n!$. The identity
$\binom nk k!(n-k)!=n!$ gives
$\E(f\star g)=\E(f)\E(g)$ coefficientwise. Also
$\E(c)=\cos(2x)$ and $\E(d)=e^x-1$. Thus
$\E(a)\cos(2x)=e^x-1$. The cosine has constant coefficient one and is
a multiplicative unit, so its inverse gives (5), uniquely in $\Q[[x]]$.
By construction $n![x^n]\E(a)=a_n$.
\end{proof}

\begin{theorem}[The two A380056 conjectures]\label{thm:egf}
For every $n\ge1$, $a_n\equiv1\pmod4$. For every $n\ge0$,
\[
 a_n\equiv2\cdot2^n+3\cdot4^n\pmod5. \tag{6}
\]
Consequently the residues modulo five at index classes $0,1,2,3$ modulo
four are respectively $0,1,1,3$, and
\[
 5\mid a_n\quad\Longleftrightarrow\quad4\mid n.
\]
In particular $5\mid a_{4n}$ for every positive $n$.
\end{theorem}

\begin{proof}
All positive-index coefficients of $c$ vanish modulo four, so $c=\delta$
in binomial convolution over $\Z/4\Z$. Reduction of $a\star c=d$
therefore gives $a=d$, proving the first assertion.

For the second, work with sequences over $\F_5$ and put $E_s(n)=s^n$,
with $0^0=1$. The binomial theorem gives
\[
 E_s\star E_t=E_{s+t}.
\]
These are divided-power coefficient sequences. No factorial denominators
or ordinary exponential series over $\F_5$ are needed. Since $-4=1$ in
$\F_5$, the sequence $c$ is one at even indices and zero at odd indices,
and hence
\[
 c=3(E_1+E_4),\qquad d=E_1-E_0.
\]
For $b=2E_2+3E_4$, bilinearity gives
\[
 \begin{aligned}
 b\star c
 &=6E_3+6E_1+9E_0+9E_3\\
 &=E_1-E_0=d
 \quad\text{over }\F_5.
 \end{aligned}
\]
The recurrence in Lemma~\ref{lem:egfsource} remains uniquely solvable
over $\F_5$, so $a=b$ after reduction, proving (6). Both $2^4$ and $4^4$
equal one modulo five. The four residues follow by evaluating the four
index classes, and only class zero gives zero.
\end{proof}

The endpoints in \texttt{Round4.ExpCosine} prove the integer convolution
recurrence, its uniqueness, the factorial coefficient interpretation,
the literal quotient in (5), uniqueness among rational formal series,
both conjectures, and the stronger complete classification modulo five.
The source bridge uses Mathlib's formal cosine series rescaled by two.

\section{Rational theta sums: one proof and one correction}\label{sec:theta}

For a formal series $U$ over a commutative ring, define
\[
 Q_j(U)=\frac{U-x^j}{1-x^jU}\quad(j\ge1),\qquad
 \Sigma(U)=\sum_{j\ge1}x^{j^2}Q_j(U)^j. \tag{7}
\]
The denominator has constant coefficient one, so the quotient means
genuine multiplicative inversion. Every summand begins in degree at
least $j^2$. Thus each coefficient in (7) is a finite sum, without any
analytic convergence assumption.

The positive-index formula explicitly displayed in A380060 is
\cite{oeis060}
\[
 A=1+2\Sigma(A), \tag{8}
\]
while the corresponding formula in A380680 is \cite{oeis680}
\[
 B=1+2\Sigma(2). \tag{9}
\]
The source pages also give bilateral expressions. The formal development
uses their positive-index formulas (8) and (9), where all quotients and
coefficients are defined directly in an ordinary formal power-series
ring.

\begin{lemma}[Existence for the implicit theta equation]\label{lem:thetaexist}
Equation (8) has a unique formal solution over any commutative ring, and
its constant coefficient is one. Its construction commutes with
coefficient ring homomorphisms.
\end{lemma}

\begin{proof}
Suppose $U-V=O(x^r)$. For constant-one units $D,E$, the identity
\[
 D^{-1}-E^{-1}=D^{-1}(E-D)E^{-1}
\]
shows that inversion preserves coefficient agreement. Hence
$Q_j(U)-Q_j(V)=O(x^r)$, and the same holds for the difference of their
$j$th powers. Multiplication by $x^{j^2}$ improves this to at least
$O(x^{r+1})$. Since each coefficient sum is finite,
\[
 (1+2\Sigma(U))-(1+2\Sigma(V))=O(x^{r+1}).
\]
The right-hand side of (8) always has constant coefficient one. The
degree-by-degree construction and uniqueness argument of
Lemma~\ref{lem:construction} therefore apply starting with the constant
coefficient. Mapping the locally finite sums, unit inverses, and the
equation gives transport by uniqueness.
\end{proof}

\begin{theorem}[A380060 modulo four]\label{thm:theta060}
The integral solution of (8) satisfies
\[
 A\equiv1+2\sum_{j\ge1}x^{j^2}\pmod4.
\]
Thus for every positive index $n$, its coefficient is two modulo four
exactly when $n$ is a square, and is divisible by four otherwise.
\end{theorem}

\begin{proof}
Equation (8) first gives $A\equiv1\pmod2$. Reduction of (7) modulo two
then gives
\[
 Q_j(A)\equiv\frac{1-x^j}{1-x^j}=1\pmod2,
 \qquad
 \Sigma(A)\equiv\sum_{j\ge1}x^{j^2}\pmod2.
\]
Doubling the latter equality lifts it to an equality modulo four by
Lemma~\ref{lem:binary}. Inserting this equality into (8) proves the
whole-series congruence. Distinct positive $j$ give distinct squares,
so its coefficient support is exactly as stated.
\end{proof}

\begin{theorem}[The corrected A380680 statement]\label{thm:theta680}
The explicit integral series (9) satisfies
\[
 B\equiv1+2\sum_{j\ge1}x^{2j^2}\pmod4.
\]
Therefore, for every $n>0$, its coefficient is two modulo four if and
only if $n=2j^2$ for some positive integer $j$, and is divisible by four
at all other positive indices.
\end{theorem}

\begin{proof}
For the constant series two, reduction modulo two gives
\[
 Q_j(2)=\frac{2-x^j}{1-2x^j}\equiv x^j\pmod2.
\]
Thus the $j$th summand of $\Sigma(2)$ reduces to
$x^{j^2}(x^j)^j=x^{2j^2}$. Apply the doubling statement in
Lemma~\ref{lem:binary} and then (9). Distinct positive $j$ yield distinct
exponents $2j^2$.
\end{proof}

\begin{corollary}[The retrieved square-index conjecture is false]
For A380680, $[x]B=4$ and $[x^2]B=6$. The retrieved claim that the
residue-two support consists of squares is false, as is its separate
claim of divisibility by four at every positive nonsquare index.
\end{corollary}

\begin{proof}
In degrees one and two, only $j=1$ can contribute to (9). Since
\[
 \frac{2-x}{1-2x}=2+3x+O(x^2),
\]
we obtain $B=1+4x+6x^2+O(x^3)$. The square index one has residue zero,
while the nonsquare index two has residue two. Both failures are proved
as explicit negations of the corresponding universal statements in Lean.
\end{proof}

The source's displayed coefficients already expose these counterexamples;
the additional theorem identifies and proves the complete corrected
pattern. The two results are deliberately not counted as two proofs of
the original square-index conjecture. The definitions, locally finite
coefficient interpretation, source uniqueness for (8), congruences, and
counterexamples are in \texttt{Round4.ThetaRational}.

\section{An elementary additional observation: A392210}\label{sec:odd}

The source gives the normalized functional equation \cite{oeis210}
\[
 D=x+\left(D\!\left(\frac{x}{1-x}\right)\right)^2. \tag{10}
\]
It describes evenness at odd indices beyond the first as an apparent
pattern. This follows directly from the square on the right.

\begin{corollary}
Equation (10) has a unique normalized integral solution. Every
odd-index coefficient beyond the first is even; equivalently,
$[x^{2n-1}]D$ is even for every $n>1$.
\end{corollary}

\begin{proof}
The normalized series $M=x/(1-x)$ is defined by the actual unit inverse
of $1-x$. Lemma~\ref{lem:square}, with $T=x$ and $P=M$, constructs $D$
and proves uniqueness for (10). Over $\F_2$, the square of any formal
series has only even exponents. Therefore the only possibly nonzero
odd-index coefficient on the right of (10) is the coefficient of $x$,
which is one. All later odd-index coefficients of $D$ vanish modulo two.
\end{proof}

The module \texttt{Round4.NestedBinomial} proves this source construction
and coefficient result. This short consequence is included as an
additional verified source observation, separately from the six
conjectural assertions in Table~\ref{tab:results}.

\section{Formal verification, corroboration, and limitations}

The new proof modules use Lean 4.19.0 and the package's pinned Mathlib
revision. Each construction is connected to the appropriate source
formula by explicit theorems. The acceptance-contract modules restate
normalizations, source equations, uniqueness, and coefficient assertions
with independently written proposition types. The principal families are:

\begin{center}\small
\begin{tabular}{@{}lp{0.58\linewidth}@{}}
\toprule
Module under \texttt{Round4} & Main mathematical scope\\
\midrule
\texttt{NestedReversion} & A195194; arbitrary weights, unique tails, finite-depth stabilization, parity.\\[3pt]
\texttt{SquareComposition} & General normalized contraction $G=T+(G\circ P)^2$.\\[3pt]
\texttt{CatalanComposition} & A393172 modulo four and A393171 modulo eight, with source identifications.\\[3pt]
\texttt{ExpCosine} & A380056; integer convolution, rational EGF bridge, both congruences, exact modulo-five pattern.\\[3pt]
\texttt{ThetaRational} & A380060; the corrected A380680 pattern and two literal counterexamples.\\[3pt]
\texttt{NestedBinomial} & A392210; canonical construction and odd-index evenness.\\
\bottomrule
\end{tabular}
\end{center}

All conclusions about coefficient indices are universally quantified.
The package's reproducibility checker recompiles first-party Lean
modules in dependency order into a fresh isolated object directory,
compiles the acceptance contracts, prints the theorem statements, and
audits their transitive axioms. The permitted standard axioms are
propositional extensionality, classical choice, and quotient soundness.
No admitted proof, custom axiom, or native-decision certificate is used.
The completed release report records the exact audited scope and
compiler commands; endpoint counts are audit counts, not counts of
mathematical discoveries.

The independent Python programs use exact integers and, for the EGF
quotient, exact rational numbers. They check the displayed source
coefficients and the literal defining equations. For nested reversion,
Lagrange extraction is compared with direct coefficient reversion and
finite-depth stabilization. For A380056, the integer convolution is
compared with rational power-series division and factorial recovery.
The Catalan and theta checks evaluate their original source formulas
separately from the modular conclusions. Each machine-readable numerical
record states its finite bounds. These checks corroborate the mapping
from OEIS to the formal statements; they do not prove the unbounded
conclusions by themselves.

Kernel verification establishes the formal propositions. Interpreting a
source correctly and establishing historical priority remain distinct
questions. The explicit contracts, exact-source computations, and AI
semantic reviews help with source interpretation but do not replace
independent human review. The analytic asymptotics, nonzero real values,
and convergence questions in the source entries are outside the scope of
these formal-series proofs.

The cumulative archive retains the earlier releases and manuscripts,
their attribution, and their already stated limitations. Earlier results
are not counted again in this paper. No OEIS edit or GitHub publication
is implied by preparing this archive.

\clearpage
\begin{thebibliography}{9}

\bibitem{oeis194}
Paul D. Hanna, \emph{OEIS A195194}, entry dated 11 September 2011;
parity conjecture dated 7 April 2026.
\url{https://oeis.org/A195194}. Retrieved 24 October 2026.

\bibitem{oeis171}
Paul D. Hanna, \emph{OEIS A393171}, entry dated 4 February 2026.
\url{https://oeis.org/A393171}. Retrieved 24 October 2026.

\bibitem{oeis172}
Paul D. Hanna, \emph{OEIS A393172}, entry dated 4 February 2026.
\url{https://oeis.org/A393172}. Retrieved 24 October 2026.

\bibitem{oeis056}
Paul D. Hanna, \emph{OEIS A380056}, entry dated 28 January 2025.
\url{https://oeis.org/A380056}. Retrieved 24 October 2026.

\bibitem{oeis060}
Paul D. Hanna, \emph{OEIS A380060}, entry dated 6 February 2025.
\url{https://oeis.org/A380060}. Retrieved 24 October 2026.

\bibitem{oeis680}
Paul D. Hanna, \emph{OEIS A380680}, entry dated 6 February 2025.
\url{https://oeis.org/A380680}. Retrieved 24 October 2026.

\bibitem{oeis210}
Paul D. Hanna, \emph{OEIS A392210}, entry dated 2 February 2026.
\url{https://oeis.org/A392210}. Retrieved 24 October 2026.

\end{thebibliography}
\end{document}
